\maketitle

\tableofcontents



\section*{The programme's questions}

The programme reconstructs structure from relations, not spaces.
Point-spaces (the real line, Stone spaces, Hilbert spaces) are
outcomes of the reconstruction, not its starting material.  This
is a bet: that the relational description --- contexts, reference
frames, directed refinement --- is the right starting point, and
that the structures usually assumed (points, $\sigma$-algebras,
global states) are \emph{reconstructed} from it.

\medskip\noindent
Three central questions organize the reading:

\begin{enumerate}
  \item \textbf{What must be added to observation to get structure?}

    The persistent question across all phases:
    \begin{itemize}[nosep]
      \item Phase 1: observation $\to$ dynamics (semigroup)
      \item Phase 2--3: observation $\to$ geometry
        (curvature, metric, action)
      \item Phase 4--6: observation $\to$ probability (CE),
        realization (descent), value-definiteness (distributivity)
    \end{itemize}

    \smallskip\noindent
    The reading form: \emph{Where does someone identify the exact
    condition that forces a valuation to behave --- and what
    structure does that condition live in?}

  \item \textbf{Reconstruction without points.}

    When can observation be modeled relationally --- via contexts
    and reference frames --- rather than by presupposing an
    underlying set-theoretic structure of points?  Numerical
    identity (point-evaluation) is the degenerate case where all
    contexts contract to a single point.

    \smallskip\noindent
    Paper~I does this for the Boolean case: probability is
    reconstructed from directed refinement of finite Boolean
    algebras, without assuming a sample space.  The open question
    is whether this extends beyond the Boolean case, and if so
    what plays the role of ``directed refinement'' when the
    algebra is not distributive.

    \smallskip\noindent
    \emph{(Sharpened by Zafiris, Ch.~6 of Foundations of
    Relational Realism: the relational description is not a
    technical convenience but the starting point; point-spaces
    are what you reconstruct, not what you assume.)}

  \item \textbf{Local-to-global extension.}

    When does locally defined information (on each Boolean context)
    determine a global object?  What are the \emph{conditions on
    the site} --- not the individual valuations --- that control
    whether local data extend globally?

    \smallskip\noindent
    CE is one instance: local charges always glue finitely
    additively (Kolmogorov), but gluing to a $\sigma$-additive
    measure requires something extra --- and that ``something
    extra'' is not a sheaf condition on any known topology
    (Zafiris 2006, Biesel 2024: dictionary translation, not
    theorem).  Paper~II's EA/PR/VDR is another instance: the
    obstruction to global state extension in the non-Boolean case
    is non-distributivity rather than failure of $\sigma$-additivity.

    \smallskip\noindent
    The general form: both the Boolean and non-Boolean extension
    problems are local-to-global questions on different categories.
    What controls the answer is the structure of the site, not the
    properties of individual states.

    \smallskip\noindent
    \emph{(Sharpened by Zafiris, Ch.~6: the gluing conditions
    on overlapping Boolean reference frames are the structural
    content; the topology of the site is what does the work.)}
\end{enumerate}

% ====================================================================
% PART I: THE EXTENSION BOUNDARY
% ====================================================================

\part{The extension boundary}

\noindent\textbf{[Status 2026-06-18: Strategy~D RETIRED (\texttt{/audit full},
prior-art --- the math is true but Gaifman 1964 inhabits the cell, stronger;
``Strategy D'' was never field-open). The LONE open problem is the OML descent /
$\sigma$-essential question (non-Boolean). The Strategy~D material below is the
closed reasoning trail; see \texttt{strategy\_d\_AUDIT\_VERDICT.md}.]}

\medskip
The programme's extension boundary historically had two faces: Strategy~D
(Boolean, set-theoretic topology --- now retired) and the OML extension/descent
problem (non-Boolean, algebraic --- the surviving open thread).  Each asks what
structural conditions force honest ($\sigma$-additive) probability.

% ====================================================================
\section{Strategy D}
% ====================================================================

\subsection*{Synthesis: the $\sigma$-complete characterization pattern}

Every known \emph{characterization theorem} for when a Boolean algebra
carries a strictly positive $\sigma$-additive probability presupposes
Dedekind $\sigma$-completeness:

\begin{itemize}[nosep]
  \item Kelley (1959), Theorem~9: complete + weakly distributive + chargeable.
  \item KVP (Fremlin 391D): Dedekind $\sigma$-complete + weakly
    $(\sigma,\infty)$-distributive + chargeable.
  \item Balcar--Jech--Paz\'{a}k (Fremlin 539N, under p-ideal dichotomy):
    Dedekind complete + ccc $\Rightarrow$ Maharam algebra.
  \item Kalton--Roberts (Fremlin 392F--G): uniformly exhaustive
    submeasure + Dedekind $\sigma$-complete.
  \item All negative constructions (Argyros 1983, Souslin line under
    $\neg$SH, Farah--Veli\v{c}kovi\'{c} 2006, Talagrand 2008) also
    operate in the $\sigma$-complete row.
\end{itemize}

\noindent
The non-$\sigma$-complete regime has no analogous characterization ---
only piecewise existence results for specific construction families
(minimally generated algebras, Borodulin-Nadzieja; Theorem~9.4 in
Plebanek 2024).  This makes Strategy~D structurally novel: not
``unstudied,'' but \emph{no measure-existence characterization reaches
here}.  Confirmed across $\geq 6$ sources: Fremlin ch39, ch53, ch54;
Kelley (1959); Argyros (1983); D\v{z}amonja--Plebanek (2008).

The most concrete ZFC construction is
Todor\v{c}evi\'{c}'s example (D\v{z}amonja--Plebanek, Theorem~3.1):
a ccc, non-$\sigma$-centred, measure-free Boolean algebra of
cardinality $\mathrm{add}(\mathcal{N})$.  However, the generating
family $\{T_a\}$ forms a chain under $\subseteq^*$, making the
generated algebra an interval algebra with atoms (see the detailed
entry below).  $\mathfrak{A}$ is \textbf{not} a Strategy~D
counterexample.  It demonstrates that measure-freeness does not
require $\sigma$-completeness, but no ZFC atomless
non-$\sigma$-complete measure-free example is known.

% ------------------------------------------------------------------
\subsubsection{Surveys and orientation}
% ------------------------------------------------------------------

\subsection*{Plebanek --- ``A survey on topological properties of
$P(K)$ spaces'' (2024)}

\textbf{READ (2026-05-20).  Verdict: most current secondary source
for Strategy~D.  Confirms the exact parameter regime is open.
Identifies near-miss constructions and a potential construction
family (minimally generated Boolean algebras).}

\medskip\noindent
\textbf{Source:} arXiv:2312.14755v3 (8 Jul 2024), 33~pages.

\medskip\noindent
\textbf{Scope.}  Surveys topological properties of $P(K)$, the space
of Radon probability measures on a compact Hausdorff space~$K$,
equipped with the weak* topology inherited from $C(K)^*$.  For
metrizable~$K$, the structure of $P(K)$ is well understood ($P(K)
\cong [0,1]^\omega$ for infinite~$K$).  For nonmetrizable compacta,
even basic properties depend on additional axioms of set theory.
The survey focuses on this nonmetrizable regime.

\medskip\noindent
\textbf{Strategy~D assessment.}  Strategy~D asks: does there exist a
compact totally disconnected space~$K$ such that (i)~$K$ has no
isolated points (= non-atomic Boolean algebra), (ii)~$K$ is not
basically disconnected ($=$ non-$\sigma$-complete Boolean algebra),
and (iii)~$K$ carries no strictly positive Radon probability?

The survey does not pose this exact question but provides the
landscape in which it sits.

\medskip\noindent
\textbf{Condition table for candidate constructions.}  Each candidate
is checked against Strategy~D's three conditions: non-atomic
(\checkmark/\texttimes/?), non-$\sigma$-complete
(\checkmark/\texttimes/?), measure-free (\checkmark/\texttimes/?).

\smallskip
\begin{tabular}{l c c c l}
  \textbf{Construction} & \textbf{Non-at.} & \textbf{Non-$\sigma$-c.}
    & \textbf{Meas-free} & \textbf{Note} \\[2pt]
  Argyros (1983, CH) & \checkmark & \texttimes & \checkmark
    & basically disconnected \\
  Kunen (1981, CH) & \checkmark & \checkmark & \texttimes
    & compact L-space; carries s.p.\ normal measure by construction \\
  Fedorchuk (1975, $\diamondsuit$) & \checkmark & \checkmark & ?
    & hereditarily separable; measure status unclear \\
  Talagrand (1980) & \checkmark & \checkmark & ?
    & Corson compact; has s.p.\ measure (Cor.~5.3) \\
  Minimally generated & \checkmark & \checkmark & ?
    & Borodulin-Nadzieja; see below
\end{tabular}

\smallskip\noindent
No candidate satisfies all three conditions.  Argyros's construction
is the closest to measure-free but fails condition~(ii): it is
basically disconnected.  Kunen's compact L-space satisfies (i) and
(ii) but the survey does not determine (iii); this is the most
promising near-miss to investigate.

\medskip\noindent
\textbf{Minimally generated Boolean algebras} (Theorem~9.4,
citing Borodulin-Nadzieja).  A Boolean algebra~$A$ is \emph{minimally
generated} if it is the closure of a chain of simple (= one-generator)
extensions starting from the two-element algebra.  Such algebras are
typically non-$\sigma$-complete and non-atomic.  Their Stone spaces
have controlled topological structure amenable to combinatorial
analysis.  Theorem~9.4 characterizes when a minimally generated BA
admits a strictly positive measure of \emph{countable Maharam type}
(via an increasing sequence of regular measures on the generating
chain).

\textbf{Why this matters for Strategy~D.}  Minimally generated
algebras automatically satisfy conditions~(i) and~(ii).  If one
could construct a minimally generated algebra failing the hypothesis
of Theorem~9.4, it would be a Strategy~D candidate --- provided the
failure of countable-type s.p.\ measures extends to s.p.\ Radon
measures of all types.  Caveat: Theorem~9.4 characterizes measures of
\emph{countable type} specifically; absence of a countable-type s.p.\
measure is weaker than absence of \emph{any} s.p.\ Radon measure.

\medskip\noindent
\textbf{Baire vs.\ Radon distinction.}  The survey works throughout
with \emph{Radon} measures on~$K$: $\sigma$-additive on the Borel
$\sigma$-algebra $\text{Bor}(K)$, inner regular with respect to
compact sets.  Paper~I's stone\_measure\_exists extends a finitely
additive charge on a Boolean algebra~$C$ to a Baire measure on
$\text{St}(C)$ (the $\sigma$-algebra generated by clopens).  For
metrizable Stone spaces, Baire $=$ Borel and every finite Baire
measure is Radon, so the distinction collapses.  For nonmetrizable
compacta (all Strategy~D candidates), Baire $\subsetneq$ Borel and a
Baire measure need not extend uniquely to a Radon measure.

This distinction matters: Strategy~D asks about strictly positive
\emph{Radon} measures.  Paper~I produces \emph{Baire} measures.  The
step from Baire to Radon is automatic in metrizable spaces but
requires additional argument in the nonmetrizable regime where
Strategy~D lives.

\medskip\noindent
\textbf{Key results for Strategy~D orientation:}
\begin{itemize}[nosep]
  \item \textbf{Cor.~2.8:} $K$ carries a strictly positive Radon
    measure iff there exists a continuous surjection from~$K$ onto a
    metrizable space with a strictly positive measure.  This is the
    Radon-level analogue of the Kelley intersection-number condition.
  \item \textbf{Thm.~4.5:} under MA$+\neg$CH, every compact space of
    weight~$< \mathfrak{c}$ carrying a strictly positive measure of
    countable Maharam type is a continuous image of the \v{C}ech--Stone
    remainder $\omega^* = \beta\omega \setminus \omega$.  This gives
    structural constraints on the topology of spaces admitting such
    measures.
  \item \textbf{\S 3.2 (Haydon's problem):} ``Given $\mu \in P(K)$
    and an uncountable family of Borel sets with positive measure, is
    there a point in~$K$ belonging to uncountably many of those
    sets?''  Consistently yes (under PFA), consistently no (under CH).
    Set-theoretic sensitivity of the entire subject.
  \item \textbf{\S 7:} recent results on measure-theoretic structure
    of Stone spaces of Boolean algebras, including conditions for
    $P(K)$ to be a Valdivia or Corson compact.
\end{itemize}

\medskip\noindent
\textbf{What the survey does NOT provide:}
\begin{itemize}[nosep]
  \item No construction satisfying all three Strategy~D conditions.
  \item No proof that such a construction is impossible.
  \item No discussion of the non-$\sigma$-complete / non-basically-disconnected
    regime as a separate question.
  \item No consistency result (positive or negative) for Strategy~D's
    exact parameter regime.
\end{itemize}

\medskip\noindent
\textbf{Contact with programme.}  \emph{[Stale 2026-06-18: Strategy~D RETIRED ---
killed on prior-art (Gaifman 1964); never field-open. The positioning below is the
closed reasoning trail.]}  Positions the (former) question precisely: it sat
in the gap between Argyros's
basically-disconnected measure-free constructions and the general
theory of s.p.\ Radon measures on zero-dimensional compacta.  The
minimally generated BA lineage (Borodulin-Nadzieja, Theorem~9.4) is
the most promising construction family to investigate.  Kunen's
compact L-space is the most promising near-miss to check against
condition~(iii).

\medskip\noindent
\textbf{Next steps suggested:}
\begin{enumerate}[nosep]
  \item Read Plebanek (2023), ``Musing on Kunen's compact L-space''
    (arXiv:2012.01849, 10pp) to determine whether Kunen's construction
    satisfies condition~(iii).
  \item Read Argyros (1983) to verify exactly which conditions it
    satisfies and confirm that basic disconnectedness is essential to
    the construction.
  \item Read D\v{z}amonja--Plebanek (2008) for chain-condition
    characterizations of s.p.\ measure existence.
\end{enumerate}

% ------------------------------------------------------------------
\subsection*{Plebanek --- ``Musing on Kunen's compact L-space''
(2023)}

\textbf{READ (2026-05-20).  Verdict: definitively settles Kunen's
L-space as \texttimes\ on Strategy~D condition~(iii).  Structural
lesson: the entire L-space approach to Strategy~D is
consistency-only, never ZFC.}

\medskip\noindent
\textbf{Source:} arXiv:2012.01849 (published \emph{Topology and its
Applications} 330, 2023), 10~pages.

\medskip\noindent
\textbf{The construction.}  Under $\mathrm{cf}(\mathfrak{c}) =
\omega_1$, Plebanek constructs a compact connected Corson compact
L-space~$K$ carrying a strictly positive normal probability
measure~$\mu$.  ``Normal'' means $\mu(B) = 0$ for every Borel
set~$B$ with empty interior (Plebanek's Definition~1.1); since
$K = \operatorname{supp}(\mu)$, the measure is strictly positive.
The space is an L-space: hereditarily Lindel\"{o}f but not separable.

\medskip\noindent
\textbf{Strategy~D verdict.}  Kunen's compact L-space satisfies
conditions~(i) (no isolated points, since connected) and~(ii) (not
basically disconnected --- in fact not even zero-dimensional, since
connected).  But condition~(iii) \emph{fails decisively}: the
construction explicitly produces a strictly positive Radon probability
by design.  This is not a gap to be investigated; it is the
construction's central purpose.

\medskip\noindent
\textbf{The Juh\'{a}sz constraint.}  Under MA$+\neg$CH, no compact
L-spaces exist (Juh\'{a}sz, classical).  This means the entire
Kunen/L-space lineage of constructions is \emph{consistency-only}:
even if one could modify a Kunen-type construction to remove the
strictly positive measure (which would mean removing its raison
d'\^{e}tre), the resulting space would not exist in ZFC\@.  For
Strategy~D, this rules out attack mode~A (ZFC construction) via the
L-space route entirely.  Any ZFC Strategy~D candidate must come from
a different construction family.

\medskip\noindent
\textbf{Additional structural results:}
\begin{itemize}[nosep]
  \item \textbf{Thm.~3.4:} $K$ has no nonmetrizable zero-dimensional
    closed subspaces.  This makes $K$ structurally remote from the
    totally disconnected setting where Strategy~D lives.
  \item \textbf{Thm.~4.3:} $C(K)$ is not isomorphic to $C(L)$ for
    any zero-dimensional~$L$.  The Banach-space structure of Kunen's
    space is incompatible with zero-dimensional compacta.
\end{itemize}

\medskip\noindent
\textbf{Contact with programme.}  Closes the Kunen near-miss:
\texttimes\ on condition~(iii), with the stronger lesson that the
L-space approach is consistency-only by the Juh\'{a}sz constraint.
The condition table in the Plebanek~(2024) entry is updated
accordingly.  The remaining near-misses to investigate are Fedorchuk
($\diamondsuit$, measure status unclear), Talagrand (Corson compact,
likely has s.p.\ measure), and minimally generated BAs
(Borodulin-Nadzieja, the most promising family).

% ------------------------------------------------------------------
\subsubsection{Constructions and examples}
% ------------------------------------------------------------------

\subsection*{Kelley --- ``Measures on Boolean algebras'' (1959)}

\textbf{Why read this.}  Foundational paper on the existence of
measures on Boolean algebras.  Introduces the intersection number
and the Kelley criterion (Theorem~4), the key chargeability
characterisation.  Also the first source for the link between
$\sigma$-additivity and weak distributivity in complete algebras.

\medskip\noindent
\textbf{What to read.}  All 14 pages; the paper is short and
self-contained.

\medskip\noindent
\textbf{Key results for Strategy~D.}

\begin{itemize}[nosep]
  \item \textbf{Theorem~4 (Kelley criterion).}  A Boolean algebra
    $B$ carries a strictly positive finitely additive measure iff
    $B \setminus \{0\}$ is a countable union of classes each with
    positive intersection number.  This is chargeability =
    Fremlin~391J.  Chargeability is necessary but not sufficient
    for $\sigma$-additive measures; Strategy~D asks about the
    $\sigma$-additive level.

  \item \textbf{Theorem~9.}  A \emph{complete} Boolean algebra
    carrying a strictly positive measure admits a strictly positive
    $\sigma$-additive measure iff it is weakly ($\omega_1$,
    $\omega$)-distributive.  This is the KVP theorem for the
    complete case.  The completeness hypothesis is load-bearing
    and cannot be dropped.

  \item \textbf{Theorem~12 (Souslin line).}  Assuming the negation
    of the Souslin hypothesis, the regular-open algebra of a Souslin
    line is complete, atomless, ccc, weakly countably distributive,
    and measure-free.  Against Strategy~D conditions: satisfies
    (i)~non-atomic and (iii)~measure-free, but fails (ii) because
    it is complete (hence $\sigma$-complete).  Same row as Argyros.

  \item \textbf{Ryll-Nardzewski addendum} (end of paper).  Gives a
    $\sigma$-measure analogue of Theorem~4: a monotone
    $\sigma$-approximation condition.  Less well-known than the
    Kelley criterion; potentially relevant if one wants to
    characterise $\sigma$-additive measures directly rather than
    passing through chargeability + weak distributivity.
\end{itemize}

\medskip\noindent
\textbf{What it does \emph{not} tell you.}  Every $\sigma$-additivity
result (Theorems~8, 9, 12) assumes completeness.  Kelley does not
address the non-$\sigma$-complete regime at all.  The paper
establishes the framework (intersection number, chargeability,
weak distributivity) but leaves Strategy~D's exact parameter regime
untouched.

\medskip\noindent
\textbf{Questions to carry forward.}
\begin{enumerate}[nosep]
  \item Does the Ryll-Nardzewski $\sigma$-approximation condition
    have content in the non-$\sigma$-complete case, or does it
    degenerate?
  \item Is there a meaningful weakening of weak distributivity that
    applies to non-$\sigma$-complete algebras?
  \item The Souslin-line example (Theorem~12) is consistency-only
    ($\neg$SH).  Are there ZFC complete measure-free examples?
    (Argyros gives one under CH; Talagrand's Maharam algebra is
    ZFC but asks a different question.)
\end{enumerate}

% ------------------------------------------------------------------
\subsection*{Argyros --- ``On compact spaces without strictly
  positive measure'' (1983)}

\textbf{Why read this.}  Constructs compact Hausdorff spaces with
property~$(*)$ (a chain condition stronger than ccc) but no strictly
positive measure.  These are the nearest known constructions to a
Strategy~D counterexample.

\medskip\noindent
\textbf{What to read.}  All 19 pages, focusing on the main
construction (the Stone--\v{C}ech compactification $\beta Y$
route) and the verification of property~$(*)$.

\medskip\noindent
\textbf{Key results for Strategy~D.}

\begin{itemize}[nosep]
  \item \textbf{Main construction.}  Builds compact spaces
    $X_n = \beta Y_n$ (Stone--\v{C}ech compactification of a
    locally compact space $Y_n$) satisfying property~$(*)$ but
    carrying no strictly positive Radon measure.  This gives
    condition (iii) --- measure-free.

  \item \textbf{Condition (i).}  The spaces $X_n$ have no isolated
    points (verified in the construction), so the corresponding
    clopen algebras are atomless.  Satisfies condition~(i).

  \item \textbf{Condition (ii) --- fails.}  The clopen algebra of
    $\beta Y_n$ is basically disconnected, i.e., $\sigma$-complete
    as a Boolean algebra.  Plebanek (2024 survey, Table~1) lists
    Argyros's spaces as basically disconnected.  Thus Argyros
    satisfies (i)~+~(iii) but fails (ii).  Same row as Kelley's
    Souslin-line example: the wrong row of the hierarchy.

  \item \textbf{Set-theoretic dependence.}  The construction
    appears to use CH or weaker fragments.  The status under
    MA~+~$\neg$CH is unclear from this paper alone.
\end{itemize}

\medskip\noindent
\textbf{Methodological warning.}  The Stone--\v{C}ech
compactification route naturally produces basically disconnected
(or even extremally disconnected) spaces, because $\beta Y$ inherits
strong completeness properties.  This is a systematic obstruction:
the standard technique for building measure-free spaces also
forces $\sigma$-completeness.  Any Strategy~D construction must
avoid this route or find a way to break the basic-disconnectedness
that $\beta Y$ provides.

\medskip\noindent
\textbf{Questions to carry forward.}
\begin{enumerate}[nosep]
  \item Is there an Argyros-style construction that avoids
    Stone--\v{C}ech compactification and produces a
    non-basically-disconnected space?
  \item Can one take a quotient or subalgebra of Argyros's algebra
    that breaks $\sigma$-completeness while preserving
    measure-freeness?
  \item The pattern: every known measure-free construction
    (Argyros, Souslin line, Gaifman) is $\sigma$-complete.
    Is this a coincidence or a structural obstruction?
\end{enumerate}

% ------------------------------------------------------------------
\subsection*{D\v{z}amonja--Plebanek --- ``Strictly positive measures
  on Boolean algebras'' (2008)}

\textbf{READ (2026-05-20).  Verdict: most directly relevant paper
for Strategy~D.  Contains the Todor\v{c}evi\'{c} ZFC example
(Theorem~3.1) whose Strategy~D status reduces to two verifiable
conditions on an explicit construction.}

\medskip\noindent
\textbf{Source:} \emph{Journal of Symbolic Logic} 73(4), 1416--1432.

\medskip\noindent
\textbf{Key results for Strategy~D.}

\begin{itemize}[nosep]
  \item \textbf{\S 1 (ZFC counterexample to approximability).}
    Starting from a Simon space (UDSP), constructs a Boolean algebra
    $\mathfrak{B}$ that is approximable (Claim~1.4) but admits no
    separable strictly positive measure (Claim~1.5).  This refutes the
    conjecture that Talagrand's ``approximability'' condition suffices
    for separable s.p.\ measures.  The algebra \emph{does} carry a
    (non-separable) strictly positive measure, so this is not
    measure-free.

  \item \textbf{\S 2, Theorem~2.1 (ccc forcing).}  For every
    $\sigma$-finite cc Boolean algebra $\mathfrak{B}$, there is a ccc
    forcing making it approximable (and of size~$\leq \mathfrak{c}$).
    If $\mathfrak{B}$ is atomless, each measure in the approximating
    sequence is nonatomic.

  \item \textbf{Corollary~2.8 (MA~+~$\neg$CH).}  Every atomless ccc
    Boolean algebra of size~$< \mathfrak{c}$ carries a strictly
    positive nonatomic measure.  Consistency result \emph{against}
    Strategy~D in the small-algebra regime.

  \item \textbf{Theorem~2.9 (chain condition characterisation).}
    $\mathfrak{A}$ carries a strictly positive nonatomic measure iff
    $\mathfrak{A}^+ = \bigcup_n B_n$ where: (i)~$B_n \subseteq
    B_{n+1}$, (ii)~$\mathrm{int}(B_n) \geq 2^{-n}$, (iii)~every
    $b \in B_n$ has disjoint $b_1, b_2 \in B_{n+1}$ below~$b$.

  \item \textbf{Corollary~2.11 (MA~+~$\neg$CH).}  For atomless
    Boolean algebras of size~$< \mathfrak{c}$, ccc $\iff$ the chain
    condition of Theorem~2.9.

  \item \textbf{Theorem~3.1 (Todor\v{c}evi\'{c} example, ZFC).}
    There exists (in ZFC) a Boolean algebra $\mathfrak{A}$ of
    size $\mathrm{add}(\mathcal{N})$ that is ccc, not
    $\sigma$-centred, and carries no strictly positive measure.
    Satisfies Strategy~D condition~(iii) outright.

  \item \textbf{Theorem~3.3 (consistency, $\mathrm{cof}(\mathcal{N}_\nu)
    = \omega_1$).}  There is a Boolean algebra of cardinality~$\omega_1$
    with a strictly positive measure but no \emph{separable} strictly
    positive measure.
\end{itemize}

\medskip\noindent
\textbf{The Todor\v{c}evi\'{c} construction (Theorem~3.1, traced to
source).}

D\v{z}amonja--Plebanek extract Theorem~3.1 from Todor\v{c}evi\'{c}'s
Theorem~8.4 in ``Chain-condition methods in topology,''
\emph{Topology and Its Applications} 101 (2000), 45--82.  The
construction (now read in full):

\smallskip\noindent
Let $K$ be the Cantor set, $Z = \{x \in K : \lim |x[i]|/i = 0\}$
(zero-density elements).  By a Ramsey argument, there exist
$A \subseteq Z$ (well-ordered under $\subseteq^*$, order-type
$= \mathrm{add}(\mathcal{N})$) and $B$ (a tower with uncountable
coinitiality).  Set $T = \{(t,n) : n \in \mathbb{N},\;
t \subseteq n\}$, and define generators:
\[
  T_a = \{(s,m) \in T : a \cap m \subseteq s\}, \qquad
  T_{(t,n)} = \{(s,m) \in T : m \geq n,\; s \cap n = t\}.
\]
Let $\mathcal{B}$ be the subalgebra of $\mathcal{P}(T)/\mathrm{fin}$
generated by both kinds.  Todor\v{c}evi\'{c} proves: $\mathcal{B}$
is ccc (Knaster), not $\sigma$-centred, every ultrafilter has
countable $\pi$-character, and every quotient $\mathcal{B}/\mathcal{U}$
is an interval algebra of order-type $< \mathrm{add}(\mathcal{N})$.

D\v{z}amonja--Plebanek take the subalgebra $\mathfrak{A}$ generated
by the $T_a$'s \textbf{alone} (``generators of the first kind'').
They verify: $\mathfrak{A}$ is ccc and not $\sigma$-centred (inherited),
and Lemma~3.2 shows no strictly positive measure exists (the generating
family $\{T_a\}$ satisfies $a \leq b \Rightarrow T_a \cdot T_b = 0$
in the quotient, forcing any measure to vanish on a dense set).

\medskip\noindent
\textbf{Strategy~D verification status (2026-05-20).  RESOLVED:
$\mathfrak{A}$ has atoms.}

The generating family $\{[T_a] : a \in A\}$ forms a \emph{chain}
in $\mathcal{P}(T)/\mathrm{fin}$: since $A$ is totally ordered
under $\subseteq^*$ and $a \subseteq^* b$ implies $T_b \subseteq^*
T_a$, the map $a \mapsto [T_a]$ reverses the $\subseteq^*$-order.
A Boolean algebra generated by a single chain is an interval algebra.
Every element of $\mathfrak{A}$ is a finite union of ``intervals''
$[T_a] \setminus [T_{a'}]$ where $a \subseteq^* a'$ in $A$.

Since $A$ has order-type $\mathrm{add}(\mathcal{N})$ (an ordinal),
every non-maximal $a \in A$ has an immediate $\subseteq^*$-successor
$a'$.  Then $[T_a] \setminus [T_{a'}]$ is:
\begin{itemize}[nosep]
  \item \emph{Nonzero:} $a' \setminus a$ is infinite (since
    $a \subsetneq^* a'$), so infinitely many $(s,m) \in T$ satisfy
    $a \cap m \subseteq s$ but $a' \cap m \not\subseteq s$.
  \item \emph{An atom:} any sub-interval $[T_p] \setminus [T_q]
    \leq [T_a] \setminus [T_{a'}]$ requires
    $a \subseteq^* p \subseteq^* q \subseteq^* a'$.  Immediate
    succession forces $p = a$, $q = a'$.
\end{itemize}

\noindent
$\mathfrak{A}$ \textbf{fails Strategy~D condition~(i)}: it is not
atomless.  D\v{z}amonja--Plebanek's Theorem~3.1 does not claim
atomlessness --- they state $|\mathfrak{A}|$, ccc, not
$\sigma$-centred, and the chain/disjointness generator condition,
but atomlessness is not among the conclusions.

\medskip\noindent
\textbf{What survives.}  $\mathfrak{A}$ is the first known ZFC
example of a ccc Boolean algebra with no strictly positive measure
and cardinality $\mathrm{add}(\mathcal{N})$.  It demonstrates that
measure-freeness can arise in $\mathcal{P}(\omega)/\mathrm{fin}$
subalgebras (outside the $\beta$-compactification route that forces
$\sigma$-completeness).  But it is \emph{not} a Strategy~D
counterexample.

\medskip\noindent
\textbf{Remaining candidates.}  Todor\v{c}evi\'{c}'s full algebra
$\mathcal{B}$ (with both $T_a$ and $T_{(t,n)}$ generators) might be
atomless --- the $T_{(t,n)}$'s have a tree structure and Claim~1 of
Theorem~8.4 shows every positive element contains some $T_a \cap
T_{(t,n)}$, consistent with splitting.  But $\mathcal{B}$ lacks the
no-s.p.-measure proof: Lemma~3.2 used the chain property of the
$T_a$-only subalgebra.  Whether $\mathcal{B}$ carries a strictly
positive measure is a separate question.

\medskip\noindent
\textbf{Synthesis observations.}
\begin{itemize}[nosep]
  \item The MA~+~$\neg$CH results (Corollaries~2.8 and 2.11) show
    that under MA~+~$\neg$CH, every atomless ccc algebra of
    size~$< \mathfrak{c}$ carries a strictly positive nonatomic
    measure.  This pushes Strategy~D toward set-theoretic
    independence.

  \item Every known ZFC measure-free construction (Argyros, Souslin
    line, Gaifman, and now $\mathfrak{A}$) either is
    $\sigma$-complete or has atoms.  No ZFC atomless
    non-$\sigma$-complete measure-free example is known.

  \item The structural reason: the standard technique
    ($\beta$-compactification) forces $\sigma$-completeness;
    the $\mathcal{P}(\omega)/\mathrm{fin}$ subalgebra route
    (Todor\v{c}evi\'{c}) avoids $\sigma$-completeness but the
    chain-generated subalgebra has atoms.  Obtaining both
    atomlessness and measure-freeness simultaneously in ZFC
    remains open.
\end{itemize}

\medskip\noindent
\textbf{Remaining questions.}
\begin{enumerate}[nosep]
  \item Does the pattern ``ZFC measure-free $\Rightarrow$
    $\sigma$-complete or atomic'' admit a structural explanation,
    or is it an artifact of the available constructions?
  \item Can the MA~+~$\neg$CH positive results (Corollary~2.8) be
    strengthened to cover all atomless ccc algebras, not just those
    of size~$< \mathfrak{c}$?
  \item Does Todor\v{c}evi\'{c}'s full $\mathcal{B}$ carry a
    strictly positive measure?  If not, is it atomless?
\end{enumerate}

% ------------------------------------------------------------------
\subsubsection{Fremlin reference chapters}
% ------------------------------------------------------------------

\subsection*{Fremlin Vol.~5, Ch.~53 --- topologies and measures III}

\textbf{Source:} Fremlin, \emph{Measure Theory} Vol.~5, Ch.~53,
\S\S\,531--539 (2008).

\textbf{Why read this.}  Strategy~D asks whether a non-$\sigma$-complete
non-atomic measure-free Boolean algebra exists.  Chapter~53 is the
definitive modern treatment of the interaction between topological
properties and measure existence on compact spaces and Boolean algebras.
The question is whether any of its powerful results reach the
non-$\sigma$-complete regime.

\textbf{What the chapter covers.}

\begin{itemize}[nosep]
  \item \S531: Maharam types and Maharam-algebra structure theory.
    Operating assumption: Dedekind complete.
  \item \S532: completion regular measures.  Focus on extending measures
    from subalgebras to completions --- not directly relevant, but
    confirms the machinery assumes $\sigma$-completeness throughout.
  \item \S533: measure-compactness and measure-compact cardinals.  Key
    concept: a compact space~$K$ is \emph{measure-compact} if every
    Baire measure extends to a Radon measure.  Results characterize when
    compact spaces carry measures with prescribed support, but the
    Boolean algebras involved are $\sigma$-complete or complete.
  \item \S534: Hausdorff measures and strong measure zero.  Metric
    setting; not relevant to Strategy~D's algebraic question.
  \item \S535--538: Liftings, density topologies, pointwise compact sets
    of measurable functions.  Background infrastructure.
  \item \S539: \textbf{Key section.}  The Balcar--Jech--Paz\'{a}k
    theorem and surrounding results on the relationship between
    weak $(\sigma,\infty)$-distributivity and Maharam-algebra structure.
\end{itemize}

\textbf{\S539 in detail.}  The main results:

\begin{itemize}[nosep]
  \item 539A--539C: A Dedekind complete Boolean algebra is a Maharam
    algebra iff it is weakly $(\sigma,\infty)$-distributive and carries
    a strictly positive Maharam submeasure.
  \item 539N (Balcar--Jech--Paz\'{a}k under p-ideal dichotomy): every
    weakly $(\sigma,\infty)$-distributive Dedekind complete ccc Boolean
    algebra is a Maharam algebra.
  \item 539Qg (Farah--Veli\v{c}kovi\'{c} 2006): under
    $2^\kappa = \kappa^+$ and $\square_\kappa$, there exist Dedekind
    complete ccc algebras that are NOT Maharam algebras (i.e.,
    non-measurable despite Dedekind completeness + ccc).
\end{itemize}

\textbf{Verdict for Strategy~D:} \emph{Every} result in \S539 assumes
Dedekind $\sigma$-completeness or full Dedekind completeness.  No
theorem, example, or remark addresses Boolean algebras that are
non-$\sigma$-complete.  This is now confirmed across $\geq 5$ sources:
Fremlin ch39, Fremlin ch53, Kelley (1959), Argyros (1983),
D\v{z}amonja--Plebanek (2008) --- all measure-theoretic results on
Boolean algebras operate in the $\sigma$-complete row.  Strategy~D's
non-$\sigma$-complete regime is genuinely uncharted.

\textbf{Questions.}

\begin{enumerate}[nosep]
  \item Does weak $(\sigma,\infty)$-distributivity have a meaningful
    analogue for non-$\sigma$-complete Boolean algebras?
  \item The Farah--Veli\v{c}kovi\'{c} consistency result shows
    Dedekind-complete ccc algebras can be non-measurable.  Could a
    variant construction drop the completeness requirement?
  \item Can the Balcar--Jech--Paz\'{a}k machinery be adapted to
    ``partial $\sigma$-completeness'' (joins exist for some but not
    all countable families)?
\end{enumerate}

% ------------------------------------------------------------------
\subsection*{Fremlin Vol.~5, Ch.~54 --- real-valued-measurable
cardinals}

\textbf{Source:} Fremlin, \emph{Measure Theory} Vol.~5, Ch.~54,
\S\S\,541--544 (2008).

\textbf{Why read this.}  Strategy~D asks whether a non-$\sigma$-complete
non-atomic measure-free Boolean algebra exists.  Chapter~54 is the
definitive treatment of when measures can live on power sets and their
quotients.  The question is whether any of its results constrain the
Strategy~D parameter regime.

\textbf{What the chapter covers.}

\begin{itemize}[nosep]
  \item \S541: Saturated ideals.  An ideal $\mathcal{I}$ of
    $\mathcal{P}\kappa$ is $\kappa$-saturated and $\kappa$-additive.
    Key result: quotient $\mathcal{P}\kappa / \mathcal{I}$ is always
    Dedekind complete (541B).  Ulam's dichotomy (541P): the quotient
    algebra is either atomless with $\kappa \leq \mathfrak{c}$
    (atomlessly-measurable), or purely atomic with $\kappa >
    \mathfrak{c}$ (two-valued-measurable).
  \item \S542: Quasi-measurable cardinals.  Cardinal arithmetic
    under measurability hypotheses.
  \item \S543: The Gitik--Shelah theorem.  If $(X, \mathcal{P}X, \mu)$
    is a probability space with singletons negligible, the Maharam
    type of $\mu$ must be large ($\geq$ the additivity of $\mu$).
    Key concept: a cardinal $\lambda$ is \emph{measure-free} iff
    there is no real-valued-measurable $\kappa \leq \lambda$.
  \item \S544: Atomlessly-measurable cardinals.  Repeated integration,
    null ideals, witnessing probabilities.
\end{itemize}

\textbf{Verdict for Strategy~D:} \emph{Not directly relevant.}
The entire chapter operates on power-set algebras $\mathcal{P}X$
and their quotients $\mathcal{P}X / \mathcal{I}$, which are always
Dedekind complete (541B).  The ``measure-free'' concept (543B(e))
concerns cardinals, not general Boolean algebras.  Strategy~D asks
about specific non-$\sigma$-complete Boolean algebras, which are not
power-set quotients.

\textbf{What it confirms.}  The chapter reinforces the pattern
from ch53: the most powerful measure-existence and
measure-nonexistence results in Fremlin operate in the Dedekind
complete/$\sigma$-complete regime.  The non-$\sigma$-complete world
that Strategy~D inhabits has no representation in this chapter.

\textbf{One indirect connection.}  If a Boolean algebra $B$ is
non-$\sigma$-complete but can be embedded into a Dedekind complete
algebra $\overline{B}$ preserving some structural property, then
ch54 results about $\overline{B}$ might constrain what happens on
$B$.  But the embedding must preserve the relevant measure-theoretic
properties, and it is unclear whether such embeddings exist in the
Strategy~D regime.

% ------------------------------------------------------------------
\subsection*{Fremlin Vol.~3, Ch.~39 --- measurable algebras}

\textbf{READ (2026-05-17 and 2026-05-20).  Verdict: answers the
``what forces $\sigma$-additivity'' question precisely for the
$\sigma$-complete case.  Weak $(\sigma,\infty)$-distributivity is
the algebraic CE condition.  The submeasure/control-measure
extension (\S\S 392--394) confirms the machinery stays in the
$\sigma$-complete row.}

\medskip\noindent
\textbf{Source:} Fremlin, \emph{Measure Theory} Vol.~3, Ch.~39,
freely available at
\texttt{https://www1.essex.ac.uk/maths/people/fremlin/chap39.pdf}.

\medskip\noindent
\textbf{391D Theorem} (Kantorovich--Vulikh--Pinsker 1950).  A
Boolean algebra $\mathfrak{A}$ is measurable iff:
\begin{enumerate}[nosep]
  \item $\mathfrak{A}$ is Dedekind $\sigma$-complete,
  \item $\mathfrak{A}$ is weakly $(\sigma,\infty)$-distributive,
  \item $\mathfrak{A}$ is chargeable (admits a strictly positive
    finitely additive functional).
\end{enumerate}

\medskip\noindent
\textbf{391J Theorem} (Kelley).  $\mathfrak{A}$ is chargeable iff
$\mathfrak{A} = \{0\}$ or $\mathfrak{A} \setminus \{0\}$ is a
countable union of sets with non-zero intersection numbers.

\medskip\noindent
\textbf{Contact with programme.}  This is the algebraic answer to
Question~3 (local-to-global extension) for the Boolean case.  The
condition controlling whether local charges extend to a
$\sigma$-additive global measure is a property of the
\emph{algebra} (the site), not the charges (the sections) ---
exactly what the programme's framing predicts.  Weak
$(\sigma,\infty)$-distributivity is the algebraic face of CE.

\medskip\noindent
\textbf{The submeasure extension (\S\S 392--394).}  The KVP theorem
characterizes measurable Boolean algebras.  \S\S 392--394 probe
whether submeasure conditions can replace or supplement the KVP
conditions --- the ``control measure problem.''  This entire
investigation presupposes Dedekind $\sigma$-completeness.

Key results:
\begin{itemize}[nosep]
  \item \textbf{392F (Kalton--Roberts).}  Uniformly exhaustive
    strictly positive submeasure $\Rightarrow$ chargeable.
  \item \textbf{392G.}  Measurable iff Dedekind $\sigma$-complete
    $+$ weakly $(\sigma,\infty)$-distributive $+$ admits a strictly
    positive uniformly exhaustive submeasure.
  \item \textbf{393 (Maharam submeasures).}  A \emph{Maharam algebra}
    = Dedekind $\sigma$-complete BA with a strictly positive
    continuous-at-0 submeasure.  Every measurable algebra is Maharam.
  \item \textbf{394 (Talagrand 2008).}  There exist Maharam algebras
    that are not measurable.  The control measure problem is resolved
    negatively.
\end{itemize}

\medskip\noindent
\textbf{Relevance to Strategy~D.}  The submeasure/control-measure
landscape is the $\sigma$-complete neighbour of Strategy~D, not
Strategy~D itself.  All results assume Dedekind $\sigma$-completeness.
Strategy~D operates in the non-$\sigma$-complete regime.  The
conceptual lesson is that measurability conditions come in layers
(chargeability $<$ Maharam submeasure $<$ uniformly exhaustive $<$
measurability), and these layers separate even in the $\sigma$-complete
case.

% ====================================================================
\section{OML extension problem}
% ====================================================================

\subsection*{Synthesis: the OML obstruction landscape}

The OML extension problem asks: what forces $\sigma$-additive measure
extension on a non-Boolean OML?  The primary sources reveal a coherent
obstruction landscape:

\begin{enumerate}[nosep]
  \item \textbf{The obstruction bites at the FA level} (Pt\'{a}k--Pulmannov\'{a}
    1994): subadditive unitality forces Boolean.  Any condition strong
    enough to give full measure-like behavior destroys non-Boolean
    structure.
  \item \textbf{The worst case exists} (Pt\'{a}k 1987): exotic OMLs have
    no states at all; states cannot extend from Boolean sublogics.
    The $\sigma$-complete case is partially resolved but open in general.
  \item \textbf{Positive results require von~Neumann--type completeness}:
    Gleason ($L(\mathcal{H})$, $\dim \geq 3$), Bunce--Wright (vN
    algebras, no Type~$I_2$), Chetcuti--Dvure\v{c}enskij (lattice effect
    algebras).  All stay within vN / JBW / effect-algebra completeness.
  \item \textbf{Intermediate conditions fail}: regularity does not force
    $\sigma$-additivity on general OMPs (Navara 1992).
  \item \textbf{The question is now well-posed} (McDonald--Bimb\'{o} 2023):
    Stone-type duality for OMLs provides the target space
    $\mathcal{S}_0(\mathfrak{A})$ on which to ask whether states extend
    to $\sigma$-additive measures.  Before 2023, the OML extension
    question lacked a precise duality-theoretic formulation.
\end{enumerate}

\noindent
The gap between ``too weak'' (regularity, non-unital states) and
``too strong'' (subadditive unitality $\Rightarrow$ Boolean; vN
completeness $\Rightarrow$ positive) appears structurally robust.
Any new positive result must exploit special structure that exotic
logics lack --- Hilbert-space geometry, operator-algebraic completeness,
or a genuinely new technique (categorical, topos-theoretic, or novel
pasting).

% ------------------------------------------------------------------
\subsubsection{Obstructions}
% ------------------------------------------------------------------

\subsection*{Pt\'{a}k--Pulmannov\'{a} (1994) --- subadditive measures
and orthomodularity}

\textbf{Source:} P.\ Pt\'{a}k and S.\ Pulmannov\'{a}, ``A measure
theoretic characterization of Boolean algebras among orthomodular
lattices,'' \emph{Czechoslovak Mathematical Journal} 44(119), 1994,
pp.~743--746.

\textbf{Why read this.}  The OML extension problem asks what structural
conditions force honest ($\sigma$-additive) probability on a
non-Boolean OML\@.  This paper identifies the precise obstruction ---
and the obstruction bites already at the finitely additive level, before
$\sigma$-additivity enters the picture.

\textbf{Main result (Theorem~1).}  An orthomodular lattice~$L$ is
Boolean if and only if $L$ is \emph{unital with respect to subadditive
states}: for every $a \neq 0$ in~$L$, there exists a state $s$ on~$L$
with $s(a) = 1$ such that $s$ is subadditive (i.e.,
$s(a \vee b) \leq s(a) + s(b)$ for all $a, b$).

\textbf{Proof chain.}
\begin{enumerate}[nosep]
  \item Prop~2: a subadditive state on an OML is a \emph{valuation}
    ($s(a \vee b) + s(a \wedge b) = s(a) + s(b)$ for all $a,b$).
  \item Prop~5: a valuation $s$ with $s(a) = 1$ forces $a$ to be
    \emph{central} (i.e., compatible with every element).
  \item Theorem~1: unital w.r.t.\ subadditive states $\Rightarrow$
    every element is central $\Rightarrow$ Boolean.
\end{enumerate}

\textbf{Why this kills na\"{\i}ve OML extension.}  The chain shows that
\emph{any} condition strong enough to give full ``measure-like''
behavior (subadditivity + unitality) forces the lattice to be Boolean.
This is not a $\sigma$-additivity issue --- the obstruction is already
at the finitely additive level.  A subadditive state on a non-Boolean
OML cannot be unital, and therefore cannot fully probe the lattice's
structure.

\textbf{Prop~6 (valuationless OMLs).}  There exist OMPs (orthomodular
posets) carrying no non-trivial valuation at all --- not even a
non-unital one.  So the gap between ``some states exist'' and ``enough
states to probe structure'' is not merely quantitative; some structures
are valuationally empty.

\textbf{Consequence for the OML extension problem.}  Any approach to
forcing $\sigma$-additivity on a non-Boolean OML must either:
\begin{enumerate}[nosep]
  \item avoid requiring subadditivity (work with genuinely
    non-subadditive measures);
  \item accept non-unitality (some elements are invisible to the
    measure theory);
  \item bypass the state space entirely (topos-theoretic or
    categorical methods); or
  \item exploit special algebraic structure (as Gleason does for
    $L(\mathcal{H})$, $\dim \geq 3$).
\end{enumerate}

\textbf{Questions.}

\begin{enumerate}[nosep]
  \item Where exactly does the proof of Prop~5 use orthomodularity
    versus distributivity?  Could a weaker lattice condition suffice?
  \item Is the unital-subadditive condition equivalent to any
    known condition in the Greechie/pasting literature?
  \item Does Gleason's theorem for $L(\mathcal{H})$ evade this
    obstruction because $L(\mathcal{H})$ fails to be unital w.r.t.\
    subadditive states, or because the Hilbert-space states satisfy
    a different structural condition?
\end{enumerate}

% ------------------------------------------------------------------
\subsection*{Pt\'{a}k (1987) --- exotic logics}

\textbf{Source:} P.\ Pt\'{a}k, ``Exotic logics,''
\emph{Colloquium Mathematicum} 54(1), 1987, pp.~1--7.

\textbf{Why read this.}  An OML (``logic'') is \emph{exotic} if it has
no states at all, and \emph{nearly exotic} if it has exactly one state.
This paper constructs both and proves universal embedding theorems:
any logic embeds into an exotic one (Theorem~3.1), and any logic with
a two-valued state embeds into a nearly exotic one (Theorem~3.2),
each with prescribed Boolean center.  The Section~4 appendix extends
these to $\sigma$-complete logics.

\textbf{READ (2026-05-17).  Verdict: populates Square~A bottom-right.
Obstruction is combinatorial (Greechie pasting), not from the center.}

\medskip\noindent
\textbf{The construction.}  An \emph{exotic logic} is a
$\sigma$-orthocomplete OMP that admits finitely additive states but no
$\sigma$-additive states.  Pt\'{a}k constructs these by
\emph{Greechie pasting}: start with a collection of Boolean algebras
(blocks), identify certain atoms across blocks so that the resulting
structure is an OMP, and show that any $\sigma$-additive state would
require the atoms of one block to sum to more than~1 (by the
identifications forcing double-counting).

\medskip\noindent
\textbf{Main results.}
\begin{itemize}[nosep]
  \item \textbf{Greechie's exotic logic:} a finite OML with no states.
    Constructed from 20 three-atom and 3 four-atom Boolean blocks
    pasted so that the atom set admits two partitions into maximal
    Boolean subalgebras with 12 and 11 classes respectively ---
    contradicting state existence.
  \item \textbf{Prop~2.1:} if $L$ is exotic, then
    $L \times \{0,1\}$ is nearly exotic (exactly one state).
  \item \textbf{Prop~3.1:} any logic embeds into an exotic logic with
    trivial center.
  \item \textbf{Theorem~3.1:} for any logic $L_1$ and Boolean algebra
    $B$, there exists an exotic logic $L$ with $L_1$ as sublogic
    and $C(L) = B$.
  \item \textbf{Theorem~3.2:} if $L_1$ has a two-valued state, the
    same holds for nearly exotic $L$.
  \item \textbf{Corollary:} any Boolean algebra embeds into a nearly
    exotic logic with arbitrary center --- ruling out state-extension
    from Boolean sublogics to the whole logic.
\end{itemize}

\textbf{Section~4: $\sigma$-complete case.}  Theorem~4.1:
\begin{enumerate}[nosep]
  \item If $L_1$ is $\sigma$-complete and $B$ is a $\sigma$-algebra
    of all subsets of a set, then $L_1$ $\sigma$-completely embeds
    into an exotic logic $L$ with $C(L) = B$.  If $L_1$ also has a
    two-valued $\sigma$-additive state, the same for nearly exotic.
  \item If $L_1$ is $\sigma$-complete with all Boolean sublogics
    finite, same conclusion for arbitrary $\sigma$-Boolean $B$.
\end{enumerate}
The general $\sigma$-complete case remains open --- the construction
does not always produce $\sigma$-complete logics.  Pt\'{a}k notes
this may be related to the existence of least $\sigma$-products
(open problem in Sikorski).

\textbf{Relevance to the OML extension problem.}  This paper
establishes the \emph{worst case} for OML state theory: state-free
OMLs exist, states cannot be extended from Boolean sublogics, and
the exotic/nearly-exotic phenomena persist in the $\sigma$-complete
setting.  Combined with Pt\'{a}k--Pulmannov\'{a} (1994), the picture
is:
\begin{itemize}[nosep]
  \item subadditive unitality forces Boolean (1994);
  \item without unitality, states can be entirely absent (1987);
  \item even one-state logics resist extension from Boolean parts.
\end{itemize}
Any positive OML extension result must exploit special structure
(Hilbert-space geometry, von~Neumann completeness) that exotic logics
lack.

\textbf{Questions.}

\begin{enumerate}[nosep]
  \item Is the $\sigma$-complete exotic embedding problem (general
    case) still open, or has it been resolved since 1987?
  \item What is the minimal size of an exotic OML?  Greechie's
    original has $\sim 70$ atoms --- is there something smaller?
  \item Do the exotic embedding theorems extend to OMPs (orthomodular
    posets), or does the lattice requirement play a role?
\end{enumerate}

% ------------------------------------------------------------------
\subsection*{Navara --- ``Regularity and $\sigma$-additivity of states''
(1992)}

\textbf{READ (2026-05-17).  Verdict: counterexample kills the
conjecture that regularity forces $\sigma$-additivity.  Not an
admissibility condition.}

\medskip\noindent
\textbf{Source:} \emph{Proc.\ AMS} 115(2), 427--429.

\medskip\noindent
\textbf{The result.}  B\'{e}aver and Cook (1989) proved: on
\emph{projection lattices of von Neumann algebras}, every
$P$-regular state (= inner regular w.r.t.\ finite-dimensional
projections) is $\sigma$-additive.  Navara shows this does NOT
generalize: there exist $\sigma$-orthocomplete quantum logics
(OMPs) on which all states are $P$-regular yet some are not
$\sigma$-additive.

\medskip\noindent
\textbf{The construction.}  Navara modifies Pt\'{a}k's exotic-logic
construction: takes a Greechie pasting that admits only non-$\sigma$-additive
states, then verifies that all states on this logic are $P$-regular
(because the finite-dimensional sub-logics approximate the full logic
sufficiently well for inner regularity, but not for
$\sigma$-additivity).

\medskip\noindent
\textbf{Contact with programme:} Clarifies the bottom-right cell of
Square~A.  The OML extension problem does \emph{not} have a simple
admissibility condition analogous to regularity.  The gap between
finitely additive states and $\sigma$-additive states on OMPs is
not bridgeable by regularity alone --- the obstruction (Pt\'{a}k's
combinatorial pasting) persists even when all states are regular.
This confirms that the non-Boolean extension problem is
structurally harder than the Boolean one (where weak
$(\sigma,\infty)$-distributivity does the job via KVP~1950).

% ------------------------------------------------------------------
\subsubsection{Positive results and state-space geometry}
% ------------------------------------------------------------------

\subsection*{Shultz --- ``A characterization of state spaces of
orthomodular lattices'' (1974)}

\textbf{READ (2026-05-17).  Verdict: orthomodularity imposes NO
geometric constraint on state spaces; but $\sigma$-additive state
spaces CAN be geometrically degenerate.}

\medskip\noindent
\textbf{Source:} \emph{J.\ Combin.\ Theory~A} 17(3), 317--328.
DOI: 10.1016/0097-3165(74)90096-X.

\medskip\noindent
\textbf{Main Theorem (p.~321).}  If $C$ is a compact convex subset
of a locally convex space, then $C \cong S(L)$ (affinely
homeomorphic to the state space) of some $\sigma$-orthomodular
lattice~$L$.  If $C$ is a rational polytope, $L$ can be chosen
finite.

\medskip\noindent
\textbf{Construction.}  Uses Greechie's $G$-spaces (the same
pasting technique as Pt\'{a}k's exotic logics): points $+$ frames
with loop constraints.  Lemma~1: any $G$-space yields a
$\sigma$-OML whose state space is affinely homeomorphic to the
weight space $W(X, \mathscr{F})$.  Lemmas~2--3: forcing techniques
(adjoining points, controlling weight values) to realize arbitrary
compact convex sets as weight spaces.

\medskip\noindent
\textbf{The negative result for $\sigma$-states (p.~327).}  The
converse is FALSE for $\sigma$-additive states.  Take $L$ $=$
Borel subsets of $\mathbb{R}$ modulo Lebesgue-null sets (a
$\sigma$-Boolean algebra).  Then $S_\sigma(L)$ has no extreme
points --- by Krein--Milman, $S_\sigma(L)$ is NOT affinely
equivalent to any compact convex subset of a locally convex space.
The $\sigma$-state space can be geometrically degenerate in ways
the full state space cannot.

\medskip\noindent
\textbf{Contact with programme:} Square~A, bottom row.
\begin{enumerate}[nosep]
  \item Orthomodularity alone imposes NO constraint on state-space
    geometry --- any compact convex set arises.  The obstruction to
    $\sigma$-additivity (Pt\'{a}k) cannot be read off from convex
    geometry.
  \item $S_\sigma(L)$ is structurally different from $S(L)$: it can
    lack extreme points entirely.  This is not just ``a smaller
    convex set'' --- it can fail to be a compact convex set at all
    in the relevant sense.
  \item Greechie pasting is the common tool behind both Shultz
    (state-space universality) and Pt\'{a}k (exotic logics with no
    $\sigma$-states).  The pasting is the fundamental technique for
    the non-Boolean column.
\end{enumerate}

% ------------------------------------------------------------------
\subsection*{Bunce--Wright (1992) --- the Mackey--Gleason problem}

\textbf{Source:} L.~J.\ Bunce and J.~D.~M.\ Wright, ``The
Mackey--Gleason Problem,'' \emph{Bulletin of the AMS} 26(2), 1992,
pp.~288--293.

\textbf{Why read this.}  The Mackey--Gleason problem asks: does every
bounded finitely additive measure on the projections of a von~Neumann
algebra extend to a bounded linear functional on the algebra?  This is
the operator-algebraic version of the OML extension problem.

\textbf{Main result (Theorem).}  Let $\mathcal{A}$ be a von~Neumann
algebra with no direct summand of Type~$I_2$.  Then every bounded
finitely additive measure on $\text{Proj}(\mathcal{A})$ extends
uniquely to a bounded linear functional on $\mathcal{A}$.

\textbf{Scope and limits.}
\begin{itemize}[nosep]
  \item The extension target is a \emph{bounded linear functional on
    the von~Neumann algebra}, not a $\sigma$-additive measure on a
    dual space like $S_0(\mathcal{A})$.  These are different questions.
  \item The von~Neumann algebra structure is load-bearing: the proof
    uses the Jordan decomposition, spectral theory, and the
    Kadison--Singer paving machinery.  None of this is available for
    a general OML.
  \item The Type~$I_2$ exclusion is genuine: $M_2(\mathbb{C})$ has
    ``exotic'' measures on its projections (the $2 \times 2$ case
    is the Kochen--Specker boundary).
  \item For the specific case of $\sigma$-additive measures on
    projections (normal states), the result is Gleason's theorem
    ($\dim \geq 3$).  Bunce--Wright generalizes from $\sigma$-additive
    to bounded FA, which is the harder direction.
\end{itemize}

\textbf{What this does NOT do for the OML extension problem.}
The OML extension problem asks: given a state~$s$ on a general OML
$A$, does $s$ extend to a $\sigma$-additive measure on the
McDonald--Bimb\'{o} dual space $S_0(A)$?  Bunce--Wright addresses:
\begin{enumerate}[nosep]
  \item a different source structure (von~Neumann projections, not
    general OML);
  \item a different extension target (bounded linear functional, not
    $\sigma$-additive measure on dual space);
  \item a different direction (FA $\to$ linear functional, not
    state $\to$ $\sigma$-additive measure).
\end{enumerate}

\textbf{What it confirms.}  Positive extension results require rich
algebraic structure.  The progression Gleason $\to$ Bunce--Wright $\to$
Chetcuti--Dvure\v{c}enskij (2003) tracks increasing generality, but
all stay within the von~Neumann / JBW / lattice-effect-algebra world
where completeness and spectral theory are available.

\textbf{Questions.}

\begin{enumerate}[nosep]
  \item Is there a purely lattice-theoretic characterization of when
    Bunce--Wright-type extension works, without reference to the
    ambient $C^*$-algebra?
  \item Does the Type~$I_2$ exclusion correspond to a lattice property
    visible in the OML setting (e.g., height~$\leq 2$ blocks)?
  \item Does the resolution of the Kadison--Singer conjecture
    (Marcus--Spielman--Srivastava 2015) have downstream consequences
    for measure extension on broader classes of OMLs, or is it
    confined to the paving problem?
\end{enumerate}

% ------------------------------------------------------------------
\subsubsection{Duality framework}
% ------------------------------------------------------------------

\subsection*{McDonald--Bimb\'{o} (2023) --- Stone-type duality for OMLs}

\textbf{Source:} J.\ McDonald and K.\ Bimb\'{o}, ``Stone-type duality
for orthomodular lattices,'' manuscript, 2023.

\textbf{Why read this.}  The classical Stone duality (1936) represents
Boolean algebras as clopen algebras of Stone spaces.  This paper extends
the duality to orthomodular lattices, providing the topological dual
space~$S_0(\mathfrak{A})$ that makes the OML extension question
well-posed: ``does a state on~$\mathfrak{A}$ extend to a
$\sigma$-additive measure on~$S_0(\mathfrak{A})$?''

\textbf{Construction.}  For an OML $\mathfrak{A}$:
\begin{itemize}[nosep]
  \item $\mathcal{S}_0(\mathfrak{A})$ is the set of all \emph{filters}
    of $\mathfrak{A}$ (not just ultrafilters, as in classical Stone
    duality).
  \item Topology: generated by sets $\sigma_0(a) = \{F \in
    \mathcal{S}_0(\mathfrak{A}) : a \in F\}$ for $a \in \mathfrak{A}$.
  \item $\mathcal{A}_0(\mathcal{S}_0(\mathfrak{A}))$: the algebra of
    clopen ``orthogonality-respecting'' subsets of the dual space.
\end{itemize}

\textbf{Main results.}
\begin{itemize}[nosep]
  \item Theorem~3.9: $\mathfrak{A} \cong
    \mathcal{A}_0(\mathcal{S}_0(\mathfrak{A}))$ --- full algebraic
    recovery.
  \item Theorem~3.11: algebraic realization --- every ``OML-space''
    arises as $\mathcal{S}_0$ of some OML.
  \item The duality is functorial, extending Stone's correspondence to
    the orthomodular category.
\end{itemize}

\textbf{Key difference from Boolean Stone duality.}  In the Boolean
case, filters and ultrafilters coincide for the purpose of duality
(every filter extends to an ultrafilter, and the ultrafilter space
suffices).  In the OML case, non-maximal filters carry essential
information --- the dual space must include \emph{all} filters.  This
is because the failure of distributivity means that the ultrafilter
extension theorem fails.

\textbf{What this provides for the OML extension problem.}  The duality
makes the extension question precise: given a state $s$ on
$\mathfrak{A}$, we now have a well-defined target space
$\mathcal{S}_0(\mathfrak{A})$ on which to ask whether $s$ extends to a
$\sigma$-additive measure.  In the Boolean case,
$\mathcal{S}_0(\mathfrak{A})$ is the classical Stone space and the
answer is yes (Carath\'{e}odory / Paper~I's stone\_measure\_exists).

\textbf{What this does NOT do.}  The paper is purely algebraic/
topological --- it establishes the duality but does not address measure
theory on $\mathcal{S}_0(\mathfrak{A})$.  No states, no measures, no
extension theorems.  The measure-theoretic question is entirely open.

\textbf{Questions.}

\begin{enumerate}[nosep]
  \item What are the topological properties of
    $\mathcal{S}_0(\mathfrak{A})$ for concrete OMLs (e.g.,
    $L(\mathcal{H})$, Greechie lattices, MO$_2$)?  Is it compact?
    Hausdorff?  Zero-dimensional?
  \item Does Gleason's theorem for $L(\mathcal{H})$ correspond to a
    natural measure-extension result on $\mathcal{S}_0(L(\mathcal{H}))$,
    or do the two live in different categories?
  \item Can the ``all filters'' structure of $\mathcal{S}_0$ explain why
    OML measure extension is harder than Boolean measure extension ---
    the larger dual space makes it harder for measures to be strictly
    positive?
  \item Is there a Riesz-type representation connecting states on
    $\mathfrak{A}$ to measures on $\mathcal{S}_0(\mathfrak{A})$?
\end{enumerate}

% ====================================================================
% PART II: THE FINITE/COUNTABLE BOUNDARY
% ====================================================================

\part{The finite/countable boundary}

Where CE lives: above finitary consistency (Deligne, G\"odel
completeness), below $\omega_1$-completeness (Karp, Esp\'indola).
The boundary is precisely located from three angles:
\begin{itemize}[nosep]
  \item \textbf{Philosophical:} Howson --- Cox's axioms produce
    finite additivity only; no ``natural extension'' to
    $L_{\omega_1,\omega}$ gives $\sigma$-additivity.
  \item \textbf{Logical:} Frot --- Deligne (enough points for
    coherent topoi) $=$ finitary completeness $=$ consistency.
    $\sigma$-additivity escapes coherence.
  \item \textbf{Topos-theoretic:} Esp\'{i}ndola --- $\omega_1$-coherent
    completeness exists but Property~$T$ on the CE site collapses
    to classical conditions.  The machinery doesn't produce new
    content at this site.
\end{itemize}

% ==================================================================
\section{Howson --- ``De Finetti, Countable Additivity, Consistency
and Coherence'' (2008)}

\textbf{READ (2026-05-17).  Verdict: identifies the same gap
(finite additivity vs.\ $\sigma$-additivity) but does not
formalize it.  TERMINOLOGICAL WARNING: Howson uses ``consistency''
and ``coherence'' as synonyms for de Finetti's criterion (finite
additivity / no Dutch book).  $\sigma$-additivity sits above both.
This is opposite to the programme's usage, where consistency $=$
finite additivity and coherence $=$ the stronger condition
including $\sigma$-additivity.}

\medskip\noindent
Howson clarifies de Finetti's position on countable additivity
and shows it is internally consistent.  Key findings:

\medskip\noindent
\textbf{The mistranslation.}  De Finetti's Italian \emph{coerenza}
$=$ ``consistency'' (\emph{consistenza}).  His 1972 book uniformly
uses ``consistency.''  The English translators (Kyburg) imposed
``coherence'' against de Finetti's own usage (p.~3).  The
programme's terminology aligns with de Finetti's original, not the
mistranslation.

\medskip\noindent
\textbf{Consistency $=$ finite additivity, precisely.}  De Finetti's
criterion: $P$ is consistent iff no finite combination of bets at
the $P$-rates guarantees a loss.  Characterizes exactly the
finitely additive probability functions.  Two properties parallel
first-order logic: \emph{absoluteness} (consistency depends only
on $P[E]$, not on which algebra $E$ sits in) and
\emph{compactness} (inconsistency detected by a finite subset).
Both fail for $\sigma$-additivity (pp.~8--9).

\medskip\noindent
\textbf{De Finetti's argument against the Dutch Book for
$\sigma$-additivity (\S\S 3--4).}  The infinite fair lottery
($p_n = 0$ for all $n$) is Dutch-Bookable only if postulate~(A)
--- ``a finite sum of fair bets is fair'' --- is extended to
countable sums.  De Finetti: this extension \emph{entails}
$\sigma$-additivity, making the Dutch Book argument circular.

\medskip\noindent
\textbf{Cox $\to$ finite additivity only (\S 6, p.~19).}  Cox's
three scale-free axioms $+$ propositional logic $\to$ finitely
additive probability.  ``Cox's result does not extend to endorsing
countable additivity.  There are infinitary propositional
languages\ldots\ $L_{\omega_1,\omega}$\ldots\ but there is no
natural extension of Cox's proof which would exploit that
additional structure.''

\medskip\noindent
\textbf{The open problem (p.~17).}  ``Something is nevertheless
missing, and that is a type of completeness theorem telling us that
the rules of probability extend to finite but not countable
additivity.''  This is precisely the gap that Frot and
Esp\'{i}ndola formalize: Deligne (finitary completeness) $=$
consistency; CE requires $\omega_1$-completeness.

\medskip\noindent
\textbf{Contact with programme.}  Howson (2008) is the primary
prior-art citation for the companion note.  He identifies the gap
(finite additivity has compactness, $\sigma$-additivity doesn't)
and states the open problem (p.~17: ``missing completeness
theorem'').  The companion note makes this precise via ultraproduct
proof.  The Frot/Esp\'{i}ndola boundary is the formal version of
Howson's philosophical observation: Deligne (finitary completeness)
$=$ consistency; CE requires something beyond
($\omega_1$-completeness).  Cox's axioms are the right starting
material for the programme: they produce finite additivity from
structural principles; $\sigma$-additivity requires something extra
that is not grounded the same way.

\medskip\noindent
\textbf{Five-traditions unification: DEAD (2026-05-18).}
Abramsky's contextuality obstruction is compact (finite covers);
the $\sigma$-additivity gap is non-compact (invisible to finite
tests).  Category error.  Removing Abramsky reduces the thesis to
Howson (2008) with fancier coordinates.  No new theorem.

\medskip\noindent
\textbf{Source:} \emph{BJPS} 59(1):1--23.

% ==================================================================
\section{Frot --- ``G\"{o}del's Completeness and Deligne's Theorem''
(2013)}

\textbf{READ (2026-05-17).  Verdict: confirms the dictionary,
places CE.}

\medskip\noindent
Proves Deligne's theorem (``coherent topoi have enough points'')
is equivalent to G\"{o}del's completeness for first-order logic.
The dictionary: models $=$ points of the topos (geometric morphisms
$\mathcal{E} \to \mathbf{Set}$); consistency $=$ non-triviality;
coherent theory $=$ finitary axioms $=$ coherent topos.

\medskip\noindent
\textbf{Original question:} Is CE's failure exactly ``the
classifying topos lacks enough points''?  \textbf{Answer: Not
exactly.}  CE is not expressible in coherent (finitary) logic
(confirming \L o\'{s}), so Deligne's theorem does not apply.  CE
lives \emph{above} the coherent level --- the guarantee of enough
points stops before reaching $\sigma$-additivity.  The right
framing: CE requires a completeness theorem for the
\emph{infinitary} fragment where $\sigma$-additivity lives.

\medskip\noindent
\textbf{What connects.}  Finite additivity IS coherent ---
Biesel's result (sheaf for finite covers) is the measure-theoretic
face of Deligne.  $\sigma$-additivity requires countable
operations, escaping coherence.  Frot makes this boundary precise:
it is the boundary between finitary and infinitary logic,
equivalently between coherent and non-coherent toposes.

\medskip\noindent
\textbf{Source:} arXiv:1309.0389.

% ==================================================================
\section{Esp\'{i}ndola --- ``Infinitary Generalizations of Deligne's
Completeness Theorem'' (2017; published \emph{JSL} 85(3), 2020)}

\textbf{READ (2026-05-17).  Verdict: places CE at $\kappa =
\omega_1$; identifies Property~$T$ as the structural condition.}

\medskip\noindent
Generalizes Deligne to $\kappa$-geometric logic (conjunctions and
disjunctions of ${<}\,\kappa$ formulas, existential quantification
over ${<}\,\kappa$ variables).  Key results:
\begin{itemize}[nosep]
  \item \textbf{Theorem~2.11.}  If $\kappa$ is regular with
    $\kappa^{<\kappa} = \kappa$, $\kappa$-geometric theories of
    cardinality ${\leq}\,\kappa$ are complete w.r.t.\
    $\mathbf{Set}$-valued models.
  \item \textbf{Corollary~3.1.}  Under same hypothesis, every
    $\kappa$-separable topos has enough $\kappa$-points.
  \item \textbf{Theorem~3.4.}  If $\kappa$ is weakly compact, every
    $\kappa$-coherent topos has enough $\kappa$-points.
  \item \textbf{Corollary~3.3.}  $\kappa$-coherent theories of
    cardinality ${\leq}\,\kappa$ are complete ($=$ Karp's
    completeness for $\mathcal{L}_{\kappa,\kappa}$).
\end{itemize}

\medskip\noindent
\textbf{Property~$T$} (Definition~2.2): an exactness condition on
the category --- every proper diagram (tree with bar) is completely
proper.  Combines strong distributivity $+$ dependent choice.  At
$\kappa = \omega$, property~$T$ is automatic (Remark~2.7),
recovering classical Deligne.  At $\kappa > \omega$ it is a genuine
restriction.  Rule~$T$ (p.~1150) is the syntactic counterpart.

\medskip\noindent
\textbf{Where CE lives.}  $\sigma$-additivity uses countable
universal quantification and countable conjunction $\to$
expressible in $L_{\omega_1,\omega}$ $\to$ relevant level is
$\kappa = \omega_1$.  At this level:
\begin{itemize}[nosep]
  \item Theorem~2.11 applies under CH
    ($\omega_1^{<\omega_1} = \omega_1$).
  \item Theorem~3.4 requires $\omega_1$ weakly compact (large
    cardinal axiom).
\end{itemize}

\medskip\noindent
\textbf{Reduction check (2026-05-17).}  Property~$T$ at
$\kappa = \omega_1$ on the site of finite Boolean subalgebras
reduces to $\sigma$-completeness of the underlying algebra $+$
uniqueness of the Carath\'{e}odory extension.  Both classical,
both already assumed in any CE setup.  No gap for a theorem.
Fifth instance of the pattern: CE-adjacent seeds die because
the topos/sheaf/logical machinery, applied to the specific site
where CE lives, restates classical conditions.

\medskip\noindent
\textbf{Status.}  The \emph{coordinates} are genuine (CE lives at
$\kappa = \omega_1$ in Esp\'{i}ndola's hierarchy, above Deligne
but within Karp).  Useful for exposition (Paper~I) and for the
Caramello reading direction (specifies what her machinery would
need to do).  Does not produce a seed.  The Caramello direction
remains open --- her Morita equivalence approach is structurally
different from ``apply Property~$T$ to this site.''

\medskip\noindent
\textbf{Source:} arXiv:1709.01967.

% ==================================================================
\section{Caramello --- \emph{Theories, Sites, Toposes} (2018)}

\textbf{READ (2026-05-17).  Verdict: the bridge technique cannot
see CE.  $\sigma$-additivity is not expressible in finitary
geometric logic, so it cannot appear as a quotient in Caramello's
framework.}

\medskip\noindent
Read Chapters 1, 3--10.  The book develops the ``bridge technique'':
given a geometric theory $\mathbb{T}$ of presheaf type with
canonical Morita equivalence
\[
  \mathbf{Sh}(\text{f.p.}\mathbb{T}\text{-mod}(\mathbf{Set})^{op},
  J_{\mathrm{at}})
  \;\simeq\;
  \mathbf{Sh}(\mathcal{C}_{\mathbb{T}'}, J_{\mathbb{T}'}),
\]
transfer properties between the syntactic site
$(\mathcal{C}_{\mathbb{T}'}, J_{\mathbb{T}'})$ and the semantic
site $(\text{f.p.}\mathbb{T}\text{-mod}(\mathbf{Set})^{op},
J_{\mathrm{at}})$ to obtain new results.

\medskip\noindent
\textbf{The discriminating constraint.}  A ``quotient'' of a
geometric theory $\mathbb{T}$ means adding a \emph{geometric}
sequent --- built from finite conjunctions, small disjunctions,
existential quantification.  $\sigma$-additivity requires
\emph{countable} universal quantification and countable
conjunction: it is not a geometric axiom.  So $\sigma$-additivity
literally cannot be expressed as a quotient of any geometric theory
in Caramello's framework.  The bridge technique cannot see the
finite/countable boundary where CE lives --- for the same
structural reason that Deligne/Frot cannot: finitary
expressiveness.

\medskip\noindent
\textbf{The infinitary escape.}  At $\kappa = \omega_1$
(Esp\'{i}ndola's generalization), $\sigma$-additivity becomes
expressible in $\kappa$-geometric logic.  But we already
established that Property~$T$ at $\omega_1$ on the CE
site reduces to classical conditions.  The bridge either can't see
CE (finitary) or reduces it to known content (infinitary).

\medskip\noindent
\textbf{What the book does provide:}
\begin{itemize}[nosep]
  \item Theorem~3.2.5 (Duality): bijection between quotients of
    $\mathbb{T}$ and subtoposes of $\mathbf{Set}[\mathbb{T}]$.
  \item Chapter~4: coHeyting algebra on subtoposes,
    Booleanization, DeMorganization.
  \item Chapter~6 (\S 6.1.1): canonical Morita equivalences for
    presheaf-type theories.
  \item Chapter~10: applications --- Fra\"{i}ss\'{e}'s theorem
    topos-theoretically, Galois representations.
  \item p.~337(b): the theory of decidable Boolean algebras is of
    presheaf type; f.p.\ models $=$ finite Boolean algebras;
    homogeneous models $=$ atomless Boolean algebras.  But this is
    the \emph{algebraic} theory --- no measures, charges, or
    $\sigma$-additivity.
\end{itemize}

\medskip\noindent
\textbf{No measure-theoretic applications.}  Caramello has no
published work applying the bridge technique to measure theory,
valuations, or probability.

\medskip\noindent
\textbf{Contact with programme:} None directly.  The bridge
technique operates within finitary geometric logic, which cannot
express $\sigma$-additivity.  Sixth instance of the pattern: CE is
invisible to every finitary framework tried (Zafiris, Biesel,
Frot/Deligne, Esp\'{i}ndola-reduced, Caramello).

\medskip\noindent
\textbf{Source:} Oxford University Press, 2018.
ISBN 978-0-19-875891-4.

% ==================================================================
\section{Biesel --- ``Sheaves of Probability'' (2024)}

\textbf{READ (2026-05-17).  Verdict: confirms dictionary translation.}

\medskip\noindent
Proves measures form a sheaf for finite covers (Thm~8),
probability measures likewise (Thm~11).  Countable covers
explicitly excluded (Remark~7: uniform measures on $\{1,\ldots,n\}$
fail to glue to $\mathbb{N}$).  Section~5 states explicitly that
the results hold for merely finitely additive charges.  For
countable covers, $\sigma$-finiteness is needed --- standard
measure theory, no hidden structure.

\medskip\noindent
\textbf{Original question:} Is the step from finite to countable
covers exactly the CE boundary?  \textbf{Answer: Yes, but
tautologically.}  Finitely additive $=$ sheaf for finite covers;
$\sigma$-additive $=$ sheaf for countable covers.  This is a
restatement, not a characterization.

\medskip\noindent
\textbf{Useful residue:} Thms~8 and~11 are clean citations for
Paper~I's CE discussion.  Ref~[2] (Ross 2012, ``All roads lead to
violations of countable additivity,'' \emph{Phil.\ Studies})
relevant for Howson/de~Finetti direction.

\medskip\noindent
\textbf{Source:} arXiv:2401.01968.

% ==================================================================
\section{Vickers --- \emph{Topology via Logic} (1989)}

\textbf{READ (2026-05-17).  Verdict: provides the framework
(locales, sublocales, Stone duality) but not the application to
measures.  The question remains open --- requires later literature.}

\medskip\noindent
The book is a domain theory / logic textbook.  It develops:
\begin{itemize}[nosep]
  \item Locales (\S 5.4): points $=$ frame homomorphisms to
    $\mathbf{2}$; spatialization/localization functors.
  \item Sublocales (\S 6.2): $=$ nuclei on frames $=$ frame
    quotients.  The sublocale of ``realized points'' (principal
    ultrafilters) is well-defined in this framework.
  \item Spectral algebraic locales (\S\S 9.2--9.3): coherent
    presentations, compact coprime opens, Stone duality.
  \item Stone spaces (\S 9.5): spectral $+$ Hausdorff
    $\leftrightarrow$ Boolean algebra.  Clopens $=$ elements of the
    Boolean algebra.  Booleanization functor.  Patch topology.
  \item $\omega$-algebraicity (\S 9.4): second countability,
    $\omega$-chain completions.  The finite/countable boundary
    appears here --- finite locales are spectral algebraic
    (Prop.~9.4.1); second countable $\leftrightarrow$
    $\omega$-algebraic (Thm.~9.4.6).
\end{itemize}

\medskip\noindent
\textbf{What the book does NOT provide:}
\begin{itemize}[nosep]
  \item No theory of valuations/measures on locales.
  \item No discussion of $\sigma$-additivity or CE in
    locale-theoretic terms.
  \item No characterization of ``a valuation is supported on the
    sublocale of principal ultrafilters.''
\end{itemize}

\medskip\noindent
\textbf{Contact with programme.}  The framework is now in hand.
The open question is whether Vickers's later localic measure theory
(2011) provides the valuation-on-locale machinery needed to
formulate ``$\mu$ is supported on the sublocale of principal
ultrafilters'' as a frame-internal condition.  This would be the
next reading direction if the CE question is revisited.

\medskip\noindent
\textbf{Source:} Cambridge University Press, 1989.  ISBN 0\,521\,36062\,5.

% ==================================================================
\section{Vickers --- ``A localic theory of lower and upper
integrals'' (2008)}

\textbf{Question while reading:} Can ``a valuation is supported on
the sublocale of principal ultrafilters'' be formulated as a
frame-internal condition on a localic valuation?  Does
$\sigma$-additivity of the valuation emerge as a locale-theoretic
property (e.g., preservation of directed joins, or a continuity
condition on the valuation map)?

\medskip\noindent
\textbf{Why.}  Vickers (1989) provides the framework (locales,
sublocales, Stone duality) but not the application to measures.
This later work develops valuations on locales constructively.
The programme's question --- is CE a frame-internal property? ---
requires this machinery.  If the support condition
``$\mu(\text{sublocale of realized points}) = 1$'' can be expressed
without reference to points, it would be the non-circular
characterization of coherence.

\medskip\noindent
\textbf{Contact with programme:} CE characterization directly.
This is the actual target that Vickers (1989) pointed toward.

\medskip\noindent
\textbf{Source:} \emph{Math.\ Logic Quarterly} 54(1), 109--123.

% ====================================================================
% PART III: PAPER II ADJACENT
% ====================================================================

\part{Paper II adjacent}

% ==================================================================
\section{R\'{e}dei \& Summers --- ``Quantum Probability Theory''
(2006)}

\textbf{READ (2026-05-17).  Verdict: review paper, Paper~II
reference.}

\medskip\noindent
Reviews quantum probability in the von Neumann algebra framework.
Classical probability $=$ abelian von Neumann algebra; quantum $=$
nonabelian.  Normal states ($=$ $\sigma$-additive) characterized
by continuity.  The type I/II/III classification controls which
probability-theoretic phenomena appear.

\medskip\noindent
\textbf{Original question:} What is the obstruction to extending
states across non-Boolean joins?  \textbf{Answer:}
Noncommutativity itself.  The type classification refines this
but no single algebraic condition beyond ``non-distributive
$+$ $\dim \geq 3$'' is isolated.

\medskip\noindent
\textbf{Useful for Paper~II:} clean reference for the parallel
between classical and quantum probability in the operator algebra
framework.

\medskip\noindent
\textbf{Source:} arXiv:quant-ph/0601158.

% ==================================================================
\section{D\"{o}ring \& Isham --- ``What is a Thing?'' (2008)}

\textbf{READ (2026-05-17).  Verdict: Paper~II reference.
Daseinisation is descent-adjacent but not CE-relevant.}

\medskip\noindent
The topos is $\mathbf{Sets}^{\mathcal{V}(\mathcal{H})^{op}}$:
presheaves on the poset $\mathcal{V}(\mathcal{H})$ of commutative
subalgebras of $B(\mathcal{H})$.  The spectral presheaf $\underline{\Sigma}$
assigns to each Boolean context $V$ its Gelfand spectrum; it has
no global elements ($=$ Kochen--Specker).  Daseinisation (\S 5.2--5.3)
maps a projection $\hat{P}$ to a sub-object of $\underline{\Sigma}$
by approximating $\hat{P}$ within each Boolean context.

\medskip\noindent
\textbf{Original question:} Does daseinisation give descent from
Stone space to realization?  \textbf{Answer:} Structurally adjacent
but different.  Daseinisation goes \emph{up} (approximating
non-Boolean by Boolean); descent goes \emph{down} (Stone space to
realized points).  The failure of global sections IS Kochen--Specker,
reformulated topos-theoretically --- the non-Boolean analogue of CE
failure, but with a different obstruction mechanism
(noncommutativity vs.\ non-$\sigma$-additivity).

\medskip\noindent
\textbf{Contact with programme:} Paper~II directly.  Not
CE-relevant.

\medskip\noindent
\textbf{Source:} arXiv:0803.0417.

% ==================================================================
\section{Epperson \& Zafiris --- \emph{Foundations of Relational
Realism} (2013)}

\textbf{READ (2026-05-17).  Verdict: does not address CE.}

\medskip\noindent
Read Chapters 6, 8, 9, 10 of the book and Zafiris (2006)
\emph{J.\ Math.\ Phys.}\ 47, 092103 (the paper that puts measures
on the site).  The Grothendieck topology $J$ (epimorphic families)
does not grade covers by cardinality --- finite, countable, and
uncountable covers satisfy $J$ uniformly.  The finite/countable
boundary where CE lives is invisible.  States are finitely additive
throughout; $\sigma$-completeness is an axiom, not a sheaf
condition.  No cohomology is computed; the paper proves a
\emph{reconstruction} theorem (counit is iso), not an obstruction
theorem.  The $H^1$ unification with Abramsky--Brandenburger does
not appear.

\medskip\noindent
\textbf{Original question:} Does Zafiris's Boolean localization
give a sheaf-theoretic formulation of when local states glue to a
global $\sigma$-additive state?  \textbf{Answer: No.}

% ====================================================================
% APPENDIX: PARKED DIRECTIONS
% ====================================================================

\appendix
\part{Parked directions}

\subsubsection*{Dzhafarov--Kujala, Contextuality by Default --- PARKED}

\noindent\textbf{Primary:} E.~N.~Dzhafarov and J.~V.~Kujala,
``Contextuality-by-Default 2.0,'' \emph{Lecture Notes in Computer
Science} 10106 (2017), 16--32.  See also Dzhafarov and Kujala,
``Context-content systems of random variables,'' \emph{Philosophical
Transactions of the Royal Society A} 375 (2017), 20160389.

\medskip\noindent
\textbf{Status: PARKED (2026-05-19).}  Read in full (LNCS chapter,
pp.~16--32).  CbD starts from the position that random variables
measured in different contexts \emph{lack a joint distribution
primitively}.  A joint distribution (coupling) is an achievement,
not a given.  Contextuality is then the failure of a maximally
connected coupling to exist.  The framework uses context-content
systems $R_q^c$ (double-indexed random variables), multimaximal
couplings, and a quasi-coupling/total-variation measure of contextuality.

\medskip\noindent\textbf{Discriminating question (as posed):} Does
CbD's coupling-existence criterion have an analogue for the FA $\to$
$\sigma$A extension problem?

\medskip\noindent
\textbf{Verdict: the question was malformed.}  The coupling problem
in CbD is the \emph{Kolmogorov extension problem}: given consistent
finite-dimensional marginals, does a joint distribution exist?  The
FA/$\sigma$A extension problem is different: given a finitely additive
charge on a Boolean algebra, does it extend to a $\sigma$-additive
measure on the $\sigma$-completion?  These are distinct extension
problems with different obstructions.  CbD's ``refinement contexts''
are finite marginal distributions that are already $\sigma$-additive;
the coupling asks whether they paste into a global $\sigma$-additive
joint.  Our problem asks whether a single charge on a single algebra
can be promoted from finite to countable additivity --- no pasting
of marginals is involved.

Additionally, CbD contextuality is \emph{compact}: it is witnessed
by finite subsystems (a finite set of contexts suffices to detect
non-existence of a coupling).  The FA/$\sigma$A gap is
\emph{non-compact}: no finite test can distinguish a $\sigma$-additive
measure from a purely finitely additive charge (this is precisely the
non-first-order-axiomatizability established in the companion note).
The compact/non-compact mismatch is the same structural reason that
Abramsky--Brandenburger sheaf contextuality does not transfer.

The Yosida--Hewitt echo (quasi-couplings decomposing into signed and
positive parts, paralleling $\mu = \mu_c + \mu_p$) is vocabulary,
not theorem.  No new content about the FA/$\sigma$A boundary.

\subsubsection*{Ellerman, Partition Logic and Information --- PARKED}

\noindent\textbf{Primary:} D.~Ellerman, ``An Introduction to
Partition Logic,'' \emph{Logic Journal of the IGPL} 22 (2014),
94--125.  See also Ellerman, ``The Logic of Partitions,''
\emph{Review of Symbolic Logic} 3 (2010), 287--350; and
\emph{New Foundations for Information Theory}, Springer (2021).

\medskip\noindent
\textbf{Status: PARKED (2026-05-19).}  Read in full.
Ellerman builds a complete dual logic where partitions (equivalence
relations) play the role that subsets play in Boolean logic.
Distinctions (``dits'') dualize elements; the partition lattice
$\Pi(U)$ is non-distributive; each partition $\pi$ generates a
Boolean core $\mathcal{B}_\pi \cong \mathcal{P}(\pi_{ns})$.
``Logical entropy'' $h(\pi) = \sum p_i(1-p_i)$ is the normalized
dit-counting measure.

\medskip\noindent
\textbf{Verdict on discriminating question:}  The dit/bit translation
\emph{is} relabeling entropy.  Logical entropy $h(\pi) = 1 -
\sum p_i^2 = 1 - \operatorname{Coll}_\mu(\mathcal{G})$, i.e., one
minus the collision probability already used in the programme's
observational-resolution notes.  The partition refinement order
recovers the cylinder-algebra filtration.  The paper stays entirely
in the finite setting --- no measures, no charges, no limit behavior,
no extension problems.  No theorem in partition logic translates to
new content about the FA/$\sigma$A boundary.

\medskip\noindent
\textbf{OML re-scoping (rejected):}  The non-distributivity of
$\Pi(U)$ does not connect to the OML extension problem.  Partition
lattices are semimodular (geometric) lattices, not orthomodular
lattices: the equivalence-relation complement is not an
orthocomplement, and the failure of distributivity in $\Pi(U)$ is
of a different structural kind than in~$L(\mathcal{H})$ or a
general OML\@.  Wrong kind of non-distributivity; no transfer.

\subsubsection*{Walley, Imprecise Probability --- PARKED}

\noindent\textbf{Primary:} P.~Walley, \emph{Statistical Reasoning
with Imprecise Probabilities}, Chapman \& Hall (1991).  For a concise
treatment, see M.~C.~M.~Troffaes and G.~de Cooman, \emph{Lower
Previsions}, Wiley (2014).

\medskip\noindent
Walley takes as primitive a \emph{set} of probability measures (a
credal set), not a single measure.  Sharp (precise) probability is
the special case where the credal set is a singleton.  The framework
has a natural hierarchy: upper/lower previsions $\to$ coherent
previsions $\to$ precise previsions $\to$ $\sigma$-additive previsions.
Each inclusion is strict.  The theory of ``natural extension'' ---
the smallest coherent prevision dominating a given assessment --- is
formally analogous to the problem of extending a finitely additive
charge.

\medskip\noindent\textbf{Assessment (read \S\S6.8--6.9, Appendix~K):
PARKED --- discriminating question was malformed.}

The discriminating question conflated two distinct axes in Walley's
hierarchy.  In his framework, the credal set
$\mathcal{M}(\underline{P})$ is the set of linear previsions
(finitely additive probabilities) dominating a lower prevision
$\underline{P}$.  Singleton $\mathcal{M}(\underline{P})$ means
\emph{precision} (sharp probability), not $\sigma$-additivity.
$\sigma$-additivity corresponds to \emph{conglomerability}: a
prevision~$P$ is $\mathcal{B}$-conglomerable if conditional and
unconditional previsions cohere across a partition~$\mathcal{B}$.
Full conglomerability (for all partitions) is equivalent to countable
additivity for precise previsions (Theorem~6.9.2, crediting
Schervish--Seidenfeld--Kadane 1984).

The framework thus has \emph{two} independent axes: (i)~precision
(sharp~vs.\ imprecise), (ii)~conglomerability ($\sigma$-additive~vs.\
merely finitely additive).  The discriminating question asked about
axis~(i) as if it were axis~(ii).

The load-bearing mathematics is SSK~(1984): conglomerability
$\leftrightarrow$ countable additivity, already in the programme via
Howson~(2008) and the companion note.  Walley's original contribution
is normative (decision-theoretic interpretation of the FA/$\sigma$A
hierarchy as a rationality axiom, his~P8), not mathematical (no new
extension theorem).  The ``natural extension'' machinery~(\S3.4) is
the lower-envelope construction, formally dual to the Yosida--Hewitt
decomposition but not producing new structural results about when
extensions are $\sigma$-additive.

\emph{W-coherence} (Appendix~K), a weakening of coherence that drops
the conglomerative principle, always admits W-coherent extensions to
larger domains.  This is a rationality-axiom distinction, not a
structural one: it does not address whether a given finitely additive
charge has a $\sigma$-additive extension.

\textbf{Cluster:} Same as Howson~(2008).  The imprecise-probability
framework redescribes the FA/$\sigma$A hierarchy in
decision-theoretic language without adding new theorems about when
$\sigma$-additive extension exists.  Vocabulary, not mathematics.

\subsubsection*{Abramsky--Brandenburger, Sheaf-Theoretic Formalism --- PARKED}

\noindent\textbf{Primary:} S.~Abramsky and A.~Brandenburger, ``The
Sheaf-Theoretic Structure of Non-Locality and Contextuality,''
\emph{New Journal of Physics} 13 (2011), 113036.

\medskip\noindent
\textbf{Re-scoped.}  The contextuality application was previously
killed: the contextuality obstruction is compact (witnessed by finite
covers), while the $\sigma$-additivity gap is non-compact (invisible
to any finite test).  These are different phenomena.

What survives is the \emph{formalism}: a presheaf of local sections
(compatible local data), a sheaf condition (existence of a global
section), and a cohomological obstruction ($H^1$ or higher) measuring
the failure.  The question is whether this formal apparatus can be
applied to the Yosida-Hewitt decomposition, where the ``local
sections'' are the restrictions of a finitely additive charge to
finite subalgebras and the ``global section'' is a $\sigma$-additive
extension.

\medskip\noindent\textbf{PARKED (2026-05-19).}  Read in full.  The
formalism does not apply, and the kill is stronger than the
compact/non-compact mismatch alone.

\textbf{The cohomology is trivially zero on the natural target.}
Stone spaces are totally disconnected compact Hausdorff; their
\v{C}ech cohomology with constant coefficients vanishes in all
positive degrees.  Applied to $\operatorname{St}(B)$ for a Boolean
algebra~$B$, the AB obstruction theory gives $H^1 = 0$
unconditionally.  The purely finitely additive component~$\mu_p$
\emph{cannot} appear as a cohomological obstruction because there is
no cohomology to obstruct.

\textbf{The underlying mismatch.}  The AB framework is
combinatorial: finite measurement covers (anti-chain families on a
finite set of observables), finite incidence matrices, LP duality for
signed measures (Theorem~5.5/5.9).  \v{C}ech cohomology in the AB
sense operates on finite simplicial complexes built from finite
measurement scenarios.  The signed-measure result (every no-signaling
model has a signed global section) is LP duality on finite matrices,
not the Yosida-Hewitt decomposition.  Extending the formalism to
countable limits would require pro-sheaf or ind-sheaf machinery that
does not exist in the AB literature and would constitute new
construction, not application.

\textbf{Verdict:} The compact/non-compact mismatch (already
identified when killing the ``five traditions'' thread) is confirmed,
and the stronger cohomological vanishing makes rescue impossible
without rebuilding the formalism.  No programme contact.

\subsubsection*{Sambin, Positive Topology --- PARKED}

\noindent\textbf{Primary:} G.~Sambin, ``Some points in formal
topology,'' \emph{Theoretical Computer Science} 305 (2003), 347--408.
See also Sambin, \emph{Positive Topology: Completing the Foundational
Trilogy} (book in preparation, drafts circulating since 2020).

\medskip\noindent
\textbf{Status: PARKED (2026-05-19).}  Read the 2003 TCS
paper in full.  The framework is purely topological/lattice-theoretic
with no measure-theoretic content.

\textbf{What Sambin provides.}  A basic pair $(X, \Vdash, S)$
relates points and observables through a forcing relation.  A formal
topology $(S, A, \triangleleft, \kappa)$ has: a cover
$\triangleleft$ (inductive, from an axiom set~$A$) playing the role
of topology; and a positivity predicate~$\kappa$ (co-inductive, from
a compatibility set~$J$) replacing the classical complement.  Formal
points are derived: a filter~$\alpha \subseteq S$ must be inhabited,
convergent, split covers, and enter~$\kappa$.

\textbf{Why the primitives are the wrong shape.}  The positivity
predicate is co-inductive --- it is the constructive replacement for
topological inhabitedness (an open is ``positive'' if it contains a
formal point).  This is \emph{not} a $\sigma$-additivity test.  The
Yosida-Hewitt decomposition requires $\sigma$-completion machinery
(countable limits, continuity from above) that lives outside formal
topology.  Even if one added measures to Sambin's framework, the
positivity predicate would not characterize the $\sigma$-additive
component: $\kappa$ detects topological non-emptiness, not
measure-theoretic continuity at countable joins.  The ``completion''
in Sambin's sense is spatial completion (adding formal points to make
a locale spatial), not measure extension.

\textbf{Verdict:} Redescribes Stone duality constructively without
engaging measure theory.  The structural disconnect is not just ``no
measures present'' but ``the available primitives (cover +
positivity) are the wrong type to express $\sigma$-additivity.''  No
programme contact.

\subsubsection*{Gisin (2021) --- PARKED}

\noindent\textbf{Primary:} N.~Gisin, ``Indeterminism in
Physics, Classical Chaos and Bohmian Mechanics,'' \emph{Erkenntnis}
86 (2021), 1553--1590.

\medskip\noindent
\textbf{Status: PARKED (2026-05-19).}  Read in full.  Physics/philosophy
paper with zero measure-theoretic content.  Proposes
``finite-information numbers'' (determined prefix + undetermined
tail); shows classical mechanics becomes non-deterministic but
empirically equivalent.  No probability, no charges, no Boolean
algebras, no extension problems.  The philosophical resonance
(infinite precision is unphysical / $\sigma$-additivity requires
infinite verification) is real but shallow: Gisin's
finite-information numbers are a physical proposal, not a
constructive real-number system in the Bishop/Brouwer sense, and
the paper develops no measure theory on such objects.  No formal
contact with the FA/$\sigma$A question.

\subsubsection*{Linnebo--Shapiro (2019) --- PARKED}

\noindent\textbf{Primary:} {\O}.~Linnebo and S.~Shapiro,
``Actual and Potential Infinity,'' \emph{No\^{u}s} 53 (2019),
160--191.

\medskip\noindent
\textbf{Status: PARKED (2026-05-19).}  Read in full.  The paper develops
modal explication of potential infinity via S4.2: liberal
potentialism preserves classical logic (potentialist translation
$\forall \to \Box\forall$, $\exists \to \Diamond\exists$);
strict potentialism yields intuitionistic logic (realizability
interpretation, Heyting arithmetic).  This concerns
quantification over potentially infinite domains, not measure
extension.  $\sigma$-additivity is a property of charges on
algebras, not a quantifier pattern; the FA/$\sigma$A gap is
already known to be non-FO-axiomatizable (Tarski); switching to
intuitionistic logic does not change which Boolean algebras admit
$\sigma$-additive measures (a set-theoretic/topological question,
not proof-theoretic).

\textbf{Note:} Hamkins's \emph{set-theoretic} potentialism
(forcing extensions of the universe) is a separate thread and
should not be conflated with Linnebo--Shapiro's
mathematical-domain potentialism.  The set-theoretic version is
directly relevant to Strategy~D, where the question ``does a
non-$\sigma$-complete non-atomic measure-free Boolean algebra
exist?''\ may have different truth values in different forcing
extensions.  If Strategy~D is pursued, Hamkins is the right
contact point --- not Linnebo--Shapiro.

% ====================================================================
\section*{How to use this document}

When reading, annotate in the margins or on a separate sheet:
\begin{enumerate}[nosep]
  \item What theorem or definition surprised you.
  \item What connected to the programme.
  \item What contradicted the knowledge map.
  \item What suggested a concrete Phase~1 seed note.
\end{enumerate}

\noindent
When a reading direction produces a concrete claim worth auditing,
extract it to \texttt{notes/unsorted/} and run Phase~2.

